\subsection{Operator-learning comparisons}
\label{sec:suite}
A common training protocol was applied to the other
problems of the operator-learning suite that \citet{batlle2024kernel}
assembled from \citet{dehoop2022cost} and \citet{lu2022comprehensive}: a
Mat\'ern kernel at the median length scale with the fit capped at 6000
points, a residual MLP mean, outputs reduced by PCA, and the stage selected on
validation. Table~\ref{tab:suite} lists the results next to the reference
values collected in that study. These comparisons use different tuning
budgets and, in some cases, different data configurations. The published
methods were tuned separately for each problem. The present results are
single runs except for Helmholtz and Navier--Stokes, which use three seeds,
exact kernel solves on 19{,}000 training pairs, and output rank 400.
Rank 40 gives PCA reconstruction errors of 1.7\% and 8.2\% on these two
problems; rank 400 reduces both below 0.01\%, well below the prediction
errors. Burgers and Darcy use 800 and 824 training pairs, respectively,
slightly fewer than their baselines, and Darcy also uses a different grid.
No result is reported for Advection~I.

The kernel stage is selected on Helmholtz, Navier--Stokes, and Darcy.
Its errors range from 1.2 times the tuned reference error on Helmholtz to
2.3 times that on Navier--Stokes. This protocol does not use the
reference kernels' Cholesky preconditioning of the input. On Burgers, the
corrected mean has an error of 2.38\%, compared with 1.93\% for the
published Fourier neural operator. On Helmholtz the
per-coordinate selection improves on the kernel alone ($1.17\%$ against
$1.22\%$) by assigning a small subset of the 400 output coordinates to the
corrected mean; on Navier--Stokes it selects the kernel almost everywhere.

\begin{table}[t]
\centering
\caption{Additional benchmark survey. Published
reference mean relative $L^2$ test error on each problem (per-problem
methods tuned by their authors), compared with the common training protocol on
training splits and grids that differ from the baselines' in the cases noted
in the text: single runs, except the two rows marked three seeds (mean
$\pm$ standard deviation, exact solve on 19{,}000 pairs, output rank 400).}
\label{tab:suite}
\begin{adjustbox}{max width=\linewidth}
\begin{tabular}{llcc}
\toprule
Problem & Map & Published reference & This work \\
\midrule
Burgers & initial condition $\to$ solution & 1.93\% (FNO) & 2.38\% \\
Darcy flow & permeability $\to$ pressure & 2.32\% (POD-DeepONet) & 2.01\% \\
Advection~I & smooth pulse $\to$ solution & $\approx0\%$ (kernel) & --- \\
Advection~II & discontinuous pulse $\to$ solution & 11.28\% (linear kernel) & 11.29\% \\
Helmholtz & wavespeed $\to$ field & 1.00\% (kernel) & 1.17\% $\pm$ 0.02 (three seeds) \\
Navier--Stokes & vorticity $\to$ vorticity & 0.12\% (kernel) & 0.28\% $\pm$ 0.00 (three seeds) \\
\bottomrule
\end{tabular}
\end{adjustbox}
\end{table}

\subsection{Empirical reflection symmetry}
\label{sec:symcheck}
The continuous problem is invariant under $x_1\mapsto1-x_1$: the domain and
boundary partition are symmetric, and the input law has a symmetric
covariance. On the grid, reflecting the load ($S$) should reflect the stress
field along the first grid coordinate ($T$), i.e.\ $G(Su)=TG(u)$. For each of the first 200 samples we
searched all 40000 loads for the nearest neighbor of the reflected load $Su_i$
and compared the corresponding outputs. For the five best-matching pairs, the
input mismatch $\|Su_i-u_j\|/\|u_j\|$ ranges over $0.27$--$0.32$, while the
output mismatch $\|Tv_i-v_j\|/\|v_j\|$ ranges over $0.085$--$0.14$; reflecting
along the second coordinate instead gives output mismatches of order one.
The smaller mismatch along the first coordinate supports that choice of
output reflection axis. The empirical mean and covariance of the training loads are
invariant under reflection to within 1.1\% and 1.5\% respectively. These
moment checks are consistent with, but do not establish, the exact input-law
invariance assumed by Proposition~\ref{prop:sym}. Reflection
averaging at test time improves twenty-nine of the thirty runs of the two
MLP variants and the refiner across ten seeds
(Section~\ref{sec:experiments}). For the full-map average,
Proposition~\ref{prop:sym} controls expected loss under its exact assumptions;
a finite test sample can show an increase.
For the ten refiner runs, identifying the supplied-field average with the
full-map average requires the channel-equivariance condition in
Supplement~\ref{sec:theory-sym}.
